\begin{lemma}
\label{editorial-source-lemma-separably-closed-connected}
Let $k$ be a separably closed field.
Let $R$, $S$ be $k$-algebras. If $\Spec(R)$ and
$\Spec(S)$ are connected, then so is
$\Spec(R \otimes_k S)$.
\end{lemma}

\begin{proof}
The source convention for connectedness
includes nonemptiness (Topology,
Definition \ref{topology-definition-connected-components}).
Thus the original $R$ and $S$ are
nonzero, and Lemma
\ref{algebra-lemma-characterize-spec-connected}
applies to them. Their tensor is
also nonzero, since extending $1$
in each factor to a $k$-basis
shows $1\otimes1\neq0$.

Suppose this tensor had a nontrivial
idempotent $e$. A finite tensor
expression for this original $e$
uses coefficients in finite type
$k$-subalgebras $R_0\subseteq R$
and $S_0\subseteq S$. These are
nonzero and have no nontrivial
idempotents, since their inclusions
in the original factors preserve
idempotents and reflect $0$ and $1$.
They therefore have connected spectra.
The map $R_0\otimes_k S_0\to R\otimes_k S$
is injective by applying the two
vector-space injections successively.
It follows that the same $e$ is
a nontrivial idempotent in this
finite type tensor: its equation
$e^2-e=0$ is reflected by that
injection. Thus it suffices to
prove connectedness for the actual
finite type subalgebras, with
the original $e$ and all its
coefficients retained.

For the finite type assertion,
use $R,S$ for the two algebras
under consideration. They are
Noetherian and have finite
nonzero numbers $n,m$ of
irreducible components by Lemma
\ref{editorial-source-lemma-Noetherian-irreducible-components}.
We induct on $n+m$. For $n=m=1$,
Lemma \ref{editorial-source-lemma-separably-closed-irreducible}
gives irreducibility, hence
connectedness, of the tensor.
If $n=1$ and $m>1$, interchange
the factors through the inverse
maps $r\otimes s\mapsto s\otimes r$
and $s\otimes r\mapsto r\otimes s$.
It remains to handle $n>1$.

Here are the component-selection
details of Topology, Lemma
\ref{topology-lemma-remove-irreducible-connected}.
Form the finite graph whose
vertices are the original
irreducible components and whose
edges record nonempty intersection.
It is connected: otherwise the
unions over two parts with no
edges between them would be
disjoint nonempty closed sets
covering the spectrum, a separation.
Choose a spanning tree, obtained
by adding adjacent new vertices
until every vertex is reached.
A leaf exists by taking an
endpoint of a longest simple
path in this finite tree.
Remove that leaf. The remaining
tree is connected, so the union
of its connected vertex spaces
is connected: add them in tree
order, each meeting the existing
union. This proves the required
choice without changing any
irreducible component.

Let $T=V(\mathfrak p)$ be the
chosen component and let $T'$
be the union of the other
$n-1$ components. Take the
original ideal $I$ to be the
intersection of their minimal
primes. Then $T'=V(I)$, by
finite union of closed sets,
and its irreducible components
are those same $n-1$ components:
their defining primes are
pairwise incomparable. Both
$T$ and $T'$ are connected
and nonempty. They meet,
since otherwise these two
closed sets would separate
the original connected spectrum.
Their intersection is exactly
$V(\mathfrak p+I)$, so
$R/(\mathfrak p+I)\neq0$.

Under the original spectrum map,
the inverse images of $T$ and
$T'$ are respectively
$$
\Spec((R/\mathfrak p)\otimes_k S),\qquad
\Spec((R/I)\otimes_k S).
$$
For each of the ideals $J$
involved, the quotient comparison
sends $[r\otimes s]$ to
$[r]\otimes s$ and has inverse
$[r]\otimes s\mapsto[r\otimes s]$,
identifying $(R\otimes_k S)/J(R\otimes_k S)$
with $(R/J)\otimes_k S$.
The induction parameters for
these two inverse images are
$1+m$ and $(n-1)+m$, both
smaller than $n+m$; hence
both are connected. Their
intersection is the inverse
image of $T\cap T'$, namely
$$
\Spec((R/(\mathfrak p+I))\otimes_k S).
$$
It is nonempty because both
original tensor factors are
nonzero vector spaces and
their unit tensor is nonzero.
The two inverse images cover
the spectrum of $R\otimes_k S$.
A union of two connected sets
with nonempty intersection is
connected: a separation would
put each in one part, and
the common point forces the
same part for both.
This finishes the induction
and excludes the original
nontrivial idempotent $e$.
\end{proof}